\section{Exact results for the effective action}\label{sec9}
In the previous section we have seen that the heat kernel can be
calculated exactly for all values of $t$ if the background satisfies certain
symmetry or smoothness conditions. In this section it will be
demonstrated that for some classes of the operators the effective
action can be calculated exactly without imposing any restrictions
on behaviour of the background fields. These are the cases when the
variation of $\zeta'(0)$ with respect to the background
can be reduced to the heat kernel coefficients which are locally
computable. Then the variation is integrated to the effective action.

The most immediate example of the variations admissible for this scheme
is quantum anomalies (cf.\ sec.\ \ref{sec7}). Historically, the first
action obtained by integration of an anomaly was the chiral Wess-Zumino
action. We do not consider this action here since such a consideration
would require a lot of additional technical devices. The reader can
consult excellent original papers 
\cite{Wess:1971yu,Witten:1983tw}\footnote{Calculation of the effective
action in two dimensional QED by integration of the anomaly may
be found in \cite{Hortacsu:1979fg}.}. 
The next subsection is devoted to another well know example,
which is the Polyakov action obtained through integration of the
conformal anomaly in two dimensions.

In sec.\ \ref{s9dual} below we consider a more complicated situation when
relevant variation of the effective action leads to a linear combination
of the heat kernel coefficients of several different operators. Typically
in this case one obtains exact relations between the effective actions
rather than the effective actions themselves.
%%%%%%%
\subsection{The Polyakov action}\label{s9polyakov}
Let us consider a two-dimensional Riemannian manifold $M$
without boundary and
a scalar field $\phi$ minimally coupled to the geometry.
This means that we choose $\Phi =0$ in the action
(\ref{nonminact}) and $\Psi =0$ in the inner product
(\ref{dilaprod}). In two dimensions any metric is conformally
flat, i.e. by a suitable choice of the coordinates one may
transform it to the conformal gauge
\begin{equation}
g_{\mu\nu}=e^{2\rho (x)} \eta_{\mu\nu},\qquad
\eta_{\mu\nu}={\rm diag} (1,1).\label{cgauge}
\end{equation} 
In this gauge
\begin{equation}
\sqrt{g}R=-2\eta^{\mu\nu} \partial_\mu\partial_\nu \rho
\,.\label{confR}
\end{equation}

The effective action $W$ depends on $\rho$ only (since there is no other
external field in this problem). We substitute (\ref{ehnT}) with
$\Phi =\Psi =0$ in (\ref{convarW}) to obtain
\begin{equation}
\delta W = \frac 1{12\pi} \int_M d^2x (\delta\rho) \eta^{\mu\nu}
\partial_\mu\partial_\nu \rho \,.\label{cgcvar}
\end{equation}
This equation can be integrated with the initial condition
$W(\rho =0)=0$. 
\begin{equation}
W=\frac 1{24\pi} \int_M d^2x\,\rho \eta^{\mu\nu}
\partial_\mu\partial_\nu \rho \,.\label{caction}
\end{equation}
By returning to the covariant notations we obtain the famous
Polyakov action \cite{Polyakov:1981rd} (see also 
\cite{Luscher:1980fr}):
\begin{equation}
W=\frac 1{96\pi} \int_M d^2x \sqrt{g} R\frac 1\Delta R \,.
\label{Polaction}
\end{equation}

Note, that on a compact manifold one has to take into account contributions
from the zero modes which lead to an additional term in the Polyakov action
\cite{Dowker:1994rt}.

The problem of finding an action whose conformal variation reproduces
the conformal anomaly can be posed
in higher dimension as well \cite{Deser:1976yx} and some physical consequences
of such actions can be studied \cite{Frolov:1979tu,Frolov:1981mz}.
Complete expressions for the conformal anomaly induced effective action
in four dimensions were obtained in \cite{Riegert:1984kt,Buchbinder:1986dc}.
Conformal action in the presence of boundaries was constructed in
\cite{Dowker:1990ue}. One should keep in mind that unless the
background is conformally flat the conformal action represents only
a part of the full effective action. Applicability and limitations
of the conformal action approach are discussed in 
\cite{Erdmenger:1997yc,Balbinot:1999vg,Gusev:1999cv}.
%%%%%%
%%\subsection{The Wess--Zumino action}\label{s9WZ}
%%%%%%
\subsection{Duality symmetry of the effective action}\label{s9dual}
In this section we deal with variations of the effective action
which can be expressed through zeta functions of several operators
(in contrast to genuine anomalous transformations of sec.\ \ref{sec7}
which involve a single zeta function each).

A rather simple example can be found in the paper \cite{Kummer:1998dc}
which considered two operators of Dirac type:
\begin{equation}
\Dir =i\gamma^\mu e^{\Psi}\partial_\mu e^{-\Phi}\,,\qquad
\Dir^\dag =i\gamma^\mu e^{-\Phi}\partial_\mu e^{\Psi} \label{twoDir}
\end{equation}
on flat two-dimensional manifold without boundary
depending on two scalar functions $\Phi$ and $\Psi$. 
Let us now consider small variations
of $\Phi$ and $\Psi$. According to (\ref{varizeta}) the zeta function
of corresponding Laplacian changes as
\begin{equation}
\delta \zeta (s,\Dir\Dir^\dag )=
-2s {\rm Tr}\left[ (\Dir\Dir^\dag )^{-s} \delta\Psi -
(\Dir^\dag\Dir )^{-s} \delta \Phi \right] \,.\label{dzPP}
\end{equation}
Variation of $\zeta'(0)$ reads:
\begin{eqnarray}
\delta \zeta'(0,\Dir\Dir^\dag )&=&
-2\left[ \zeta(0,\delta \Psi,\Dir\Dir^\dag )-
\zeta (0,\delta\Phi,\Dir^\dag\Dir)\right] \nonumber\\
&=&-2\left[ a_2(\delta\Psi,\Dir\Dir^\dag)
-a_2(\delta\Phi,\Dir^\dag \Dir )\right]\,.\label{dzpPP}
\end{eqnarray} 
Here we used that $\Dir\Dir^\dag$ and $\Dir^\dag\Dir$ are operators
of Laplace type, and, therefore, the corresponding zeta functions
are regular at $s=0$. In the last line we replaced $\zeta (0)$ 
by the heat kernel coefficients by means of (\ref{anz0}).

Next we have to rewrite $\Dir\Dir^\dag$ and $\Dir^\dag\Dir$ in the
canonical form (\ref{laplaceb}) with the help of (\ref{oab}),
(\ref{Eab}). For $\Dir\Dir^\dag$ we have:
\begin{eqnarray}
&&g^{\mu\nu}=e^{2(\Psi -\Phi )}\,,\qquad \omega_\mu =
\Psi_{,\mu} -\gamma^\nu\gamma_\mu \Phi_{,\nu} \,,\nonumber\\
&&E={\Phi_{,\mu}}^\mu \,,\qquad R =-2(\Phi -\Psi {)_{,\mu}}^\mu\,,
\label{DDdef}
\end{eqnarray}
where comma denotes ordinary partial derivatives, all contractions
are performed with the effective metric $g^{\mu\nu}$. Note, that in
two dimensions $\Gamma^\rho_{\mu\nu}g^{\mu\nu}=0$ in the conformal
coordinates. To obtain relevant geometric quantities for $\Dir^\dag\Dir$
one has to replace $\Phi$ by $-\Psi$ and vice versa. Equation
(\ref{a2nobou}) gives:
\begin{equation}
\delta \zeta'(0,\Dir\Dir^\dag )=-\frac 1{6\pi}\int d^2x
\left[ \delta\Psi (\Delta \Psi +2\Delta\Phi )
+\delta\Phi (\Delta\Phi +2\Delta\Psi )\right]\,,\label{zDDvar}
\end{equation}
where $\Delta$ is the flat space Laplacian. Let us recall, that 
$\zeta'(0,D)$ defines the functional determinant of $D$ (see
eq.\ (\ref{detzeta})). The variation (\ref{zDDvar}) can be integrated
to give\footnote{The heat kernel coefficients $a_2(x,\Dir\Dir^\dag )$
and $a_2(x,\Dir^\dag\Dir )$ are total derivatives. Therefore,
the second term on the right hand side of (\ref{detzeta}) does not
contribute.}
\begin{equation}
\ln\det (\Dir\Dir^\dag)=-\zeta'(0,\Dir\Dir^\dag)=
\frac 1{12\pi} \int d^2x (\Psi \Delta \Psi +4\Psi\Delta\Phi
+\Phi\Delta\Phi )\,.\label{detDDfin}
\end{equation}
We see from (\ref{detDDfin})
that $\ln \det (\Dir\Dir^\dag)=\ln \det (\Dir^\dag\Dir )$
which is obviously true 
up to possible topological contributions from zero modes which have
been neglected in this calculation. The result (\ref{detDDfin}) was
confirmed in \cite{Gusev:1999cv} by the methods of covariant perturbation
theory.

This example is clearly of (at least an academic) 
interest since it adds up to very
few cases discussed above in this section when a closed analytic
expression for the determinant may be obtained without any assumptions 
on smoothness of the background fields.    
The operator $\Dir\Dir^\dag$ shares some similarities with 
the kinetic operator
for non-minimally coupled scalars in two dimensions (\ref{ehnA}).
There are good grounds to believe 
\cite{Kummer:1999zy} the functional determinant (\ref{detDDfin})
may describe the spherically symmetric part of the Hawking flux in four
dimensions.

Let us consider the dilatonic operator (\ref{ehnA}).
We take $\Psi =\Phi$ for simplicity\footnote{Variation of $\det D$
with respect to $\Psi$ may be expressed through the scale anomaly.
Therefore, it is not a problem to find exact dependence of $\det D$
on $\Psi$ by using the methods of sec.\ \ref{sec7}.}. We shall be
interested in properties of $\det D$ under reflection of the dilaton,
$\Phi\to -\Phi$. The quantity $\ln \det D (\Phi) - \ln \det D(-\Phi)$
defines the dilaton shift under the $T$-duality transformations
in string theory. In the string theory context this problem was
solved in \cite{Schwarz:1993te} basing on earlier results of
\cite{Schwarz:1979ae}. Here we follow presentation of 
the paper \cite{Vassilevich:2000kt}.

Let us introduce two first order differential operators
\begin{equation}
L_\mu =e^{-\Phi}\partial_\mu e^{\Phi}\,,\qquad
L_\mu^\dag =-e^{\Phi}\partial_\mu e^{-\Phi}
\label{LLdag}
\end{equation}
and associated second order operators
\begin{equation}
D_+=L_\mu^\dag L_\mu \,,\qquad 
D_-=L_\mu L_\mu^\dag \,. \label{DpDm}
\end{equation}
We restrict ourselves to flat two-dimensional manifolds.
Therefore, position of the indices (up or down) plays no role.

We calculate again variation of the $\zeta$ function with respect
to variation of $\Phi$:
\begin{equation}
\delta \zeta (s,D_+)=-2s {\rm Tr} \left[
(\delta \Phi ) \left( L_\mu^\dag L_\mu D_+^{-s-1}
-L_\mu D_+^{-s-1} L_\mu^\dag \right)\right]\,.\label{dudz}
\end{equation}
The operators in the first term here recombine in $D_+^{-s}$. The situation
with the second term is more subtle. Strictly speaking, to treat this
term one has to perform the Mellin transform and use analytic continuation
in $s$  (see footnote \ref{ft} in sec.\ \ref{s7conf}). However, the
result of such manipulations is almost obvious:
\begin{equation}
{\rm Tr} \left[ (\delta \Phi) L_\mu D_+^{-s-1} L_\mu^\dag \right]=
{\rm Tr} \left[ (\delta \Phi) H^{-s} \right] \,,\label{duvarH}
\end{equation}
where
\begin{equation}
H_{\mu\nu}=L_\mu L_\nu^\dag \label{dudefH}
\end{equation}
is a matrix operator acting on the space $V_+$ of the vector fields
which can be represented as $v_\nu = L_\nu f$ with some scalar
function $f$. The operator (\ref{dudefH}) is not of Laplace type
as it contains a complicated differential projector on $V_+$.
Therefore, even though the variation (\ref{dudz}) looks very
similar to (\ref{dzPP}), for example, it cannot be used directly
to evaluate $\zeta'(0,D_+)$ since we cannot even guarantee that
(\ref{duvarH}) is regular at $s=0$ (in fact, it is {\it not}
regular, see \cite{Vassilevich:2000kt}).  

We also have:
\begin{equation}
\delta \zeta (s,D_-)=2s {\rm Tr}\left[ (\delta\Phi) \left(
D_-^{-s} - \bar H^{-s} \right) \right] \,,\label{dudmin}
\end{equation}
where the operator
\begin{equation}
\bar H_{\mu\nu}=\epsilon_{\mu\mu'}\epsilon_{\nu\nu'} L_{\mu'}^\dag
L_{\nu'} \label{defbarH}
\end{equation}
acts on the space $V_-$ of the vector fields of the form $\bar v_{\mu}=
\epsilon_{\mu\mu'}L_{\mu'}^\dag$. All epsilon tensors cancel after
taking the trace in (\ref{dudmin}) and we arrive at the same formula
as for $D_+$ but with the replacement $\Phi\to -\Phi$. The spaces
$V_+$ and $V_-$ are orthogonal, and $V_+ + V_-=V$ is the space of
all vector fields on $M$ (we neglect possible zero modes). Therefore,
\begin{equation}
\delta \left[ \zeta (s,D_+)-\zeta (s,D_-) \right]
=2s \left[ \zeta (s,\delta\Phi,H+\bar H )-
\zeta (s,\delta\Phi,D_+) -\zeta (s,\delta\Phi,D_-) \right]\,.
\label{duzz}
\end{equation}
The operator
\begin{equation}
(H+\bar H)_{\mu\nu}=\delta_{\mu\nu} D_+ +\Phi_{,\mu\nu}
\label{HbarH}
\end{equation}
is of Laplace type. We can now act in the standard manner to
obtain
\begin{equation}
\delta \left[ \zeta' (0,D_+)-\zeta' (s,D_-) \right]
=2\left[ a_2(\delta\Phi,H+\bar H )-
a_2(\delta\Phi,D_+) -a_2(\delta\Phi,D_-) \right]\,.
\label{duzzpr}
\end{equation}
The right hand side can be evaluated by using (\ref{a2nobou}).
We leave it is an exercise to show that the variation (\ref{duzzpr})
is zero (alternatively, one may look up in \cite{Vassilevich:2000kt}).  
This leads us to the conclusion that
\begin{equation}
\zeta'(0,D_+) -\zeta'(0,D_-)=-\ln\det D_+ +\ln \det D_-=0.
\label{detdet}
\end{equation}

In the one-dimensional case a similar relation may be obtained by 
methods of supersymmetric quantum mechanics. In this simplest case,
and in higher dimensions if $\Phi$ depends on one of the coordinates only,
$D_+$ and $D_-$ are isospectral up to zero modes and (\ref{detdet})
follows immediately. In two dimensions, $D_+$ and $D_-$ are
{\it not} isospectral in general. Although $D_+$ and $D_-$ have
coinciding determinants, other spectral functions can be different.

Extensions of (\ref{detdet}) to the case of matrix-valued $\Phi$,
manifolds with boundaries, and the dilaton-Maxwell theory in
four dimensions can be found in \cite{Vassilevich:2000kt}.
A generalisation to arbitrary dimension has been achieved recently 
\cite{Gilkey:2002qd}. Let $A_p$ be a $p$-form field with the field
strength $F=dA_p$. Consider the classical action
\begin{equation}
{\mathcal{L}}=\int_M e^{-2\Phi} F\wedge \star F \,,\label{pformact}
\end{equation}
where $\star$ is the Hodge duality operator. 
Such actions appear, for example, in extended supergravities
\cite{Cremmer:1978tt} and bosonic M-theory \cite{Horowitz:2000gn}.
Instead of the dilaton, also a tachyon coupling may appear
\cite{Klebanov:1998yy}. We are interested in the symmetry properties
of the effective action under the transformation 
$p\to n-p-2$, $\Phi\to -\Phi$. In higher dimensional supergravity
theories this is a part of the $S$-duality transformation.

It is convenient to define
the twisted exterior derivative
\begin{equation}
d_\Phi :=e^{-\Phi} d e^{\Phi} \label{twider}
\end{equation}
and the associated twisted coderivative and twisted Laplacian
\begin{equation}
\delta_\Phi :=e^\Phi \delta e^{-\Phi}\,,\qquad
\Delta_\Phi :=\left( d_\Phi + \delta_\Phi \right)^2\,.\label{twicoL}
\end{equation}
Since $d_\Phi^2=0$ we have an elliptic complex. The restriction
of $\Delta_\Phi$ on the space of $p$ forms will be denoted by
$\Delta_\Phi^p$. This twisted de Rham
complex was introduced by Witten in the context of Morse theory and
supersymmetric quantum mechanics \cite{Witten:1982im}. 

Any $p$-form can be decomposed as the sum of a twisted exact,
twisted co-exact, and twisted harmonic form:
\begin{equation}
A_p=d_\Phi A_{p-1} +\delta_\Phi A_{p+1} +\gamma_p\,,
\qquad \Delta_\Phi \gamma_p=0 \,.\label{twi-decomp}
\end{equation}
The projections on the spaces of twisted exact and twisted co-exact
forms will be denoted by the subscripts $\parallel$ and $\perp$
respectively.

We assume that the fields $\tilde A_p=e^{-\Phi}A_p$ have a standard
Gaussian measure and are to be considered as fundamental fields in the
path integral. The action given in eq. (\ref{pformact}) is invariant
under the gauge transformation which sends $\tilde A_p$ to $\tilde A_p +
d_\phi \tilde A_{p-1}$. This means that the $p$-forms which are
$d_\phi$ exact have to be excluded from the path integral, but that
a Jacobian factor corresponding to the ghost fields $\tilde A_{p-1}$
has to be included in the path integral measure. Next we note that
$d_\phi$-exact $(p-1)$-forms do not generate a non-trivial transformation
of $\tilde A_p$. Hence, such fields must be excluded from the ghost
sector. Then we have to include ``ghosts for ghosts''. This goes on
until the zero forms have been reached. By giving these arguments
an exact meaning, one arrives at the Faddeev--Popov approach to
quantisation of the $p$-form actions 
\cite{Duff:1980qv,Siegel:1980jj,Obukhov:1982dt,Barbon:1995zh}. 
We note that the procedure of 
\cite{Obukhov:1982dt,Barbon:1995zh} is valid also in the
presence of a dilaton interaction if one simply replaces the ordinary
derivatives by the twisted ones. As a result, we have the following 
expression for the
effective action: 
\begin{equation}
W_p(\Phi )=\frac 12 \sum_{k=0}^p (-1)^{p+k} \ln \det (\Delta^k_\Phi
\vert_{\perp}) +W_p^{\rm top} \,.\label{twiefa}
\end{equation}
$W_p^{\rm top}$ depends on certain topological characteristics
of the manifold (the Betti numbers). We shall neglect $W_p^{\rm top}$
and some other topological contributions 
in what follows.

By combinatorial arguments, similar to the presented above,
one can show that \cite{Gilkey:2002qd}:
\begin{equation}
\delta \left( W_p(\Phi )-W_{n-p-2}(-\Phi )\right)=
\sum_{k=0}^m (-1)^{p+k}  a_n(\delta\Phi,\Delta_{\Phi}^k)\,.
\label{twipvar}
\end{equation}
One should distinguish between $\delta \Phi$ (which is a variation
of the dilaton) and $\delta_\Phi$ (which is a twisted co-derivative).
Therefore, we have related the variation of the effective actions
with respect to $\Phi$ to a combination of the heat kernel coefficients
which is called the supertrace of the twisted
de Rham complex. A somewhat surprising feature of the supertrace 
is that it can be calculated for any $n$, with or without boundaries
\cite{Gilkey:2002qd}. For example, the volume term in (\ref{twipvar})
does not depend on $\Phi$ and, therefore, is the standard
Euler density which is given by
\begin{equation}
{\rm \bf E}_n:=(4\pi)^{-\bar n}\textstyle\frac1{2^{\bar n}\bar n!}
\epsilon^{i_1...i_m}\epsilon^{j_1...j_n}R_{i_1i_2j_2j_1}...
R_{i_{n-1}i_nj_{n}j_{n-1}}
\label{Pengui}\end{equation}
(with $\bar n=n/2$) for $n$ even. ${\rm \bf E}_n=0$ for $n$ odd.

It is known \cite{Fradkin:1985ai,Barbon:1995zh} that for $\Phi =0$
the dual theories are quantum equivalent:
\begin{equation}
W_p(0)-W_{n-p-2}(0)=0 \,.\label{Phi0}
\end{equation}
By using this equation as an ``initial condition'' we can integrate
the variation (\ref{twipvar}) to obtain:
\begin{equation}
W_p(\Phi )-W_{n-p-2}(-\Phi )=(-1)^p \iM \Phi {\rm \bf E}_n \,.\label{Phidual}
\end{equation}

Further terms $\sum_p (-1)^p a_k (x,\Delta_\Phi^p)$ are of some
interest in mathematics. They can be calculated 
\cite{Gilkey:2002hk,Gilkey:2002ts} for 
$k\le n+2$ and any $n$.


